\documentclass{article}

\usepackage{neurips_2024}
\usepackage[utf8]{inputenc}
\usepackage[T1]{fontenc}
\usepackage{hyperref}
\usepackage{url}
\usepackage{booktabs}
\usepackage{amsmath}
\usepackage{amsfonts}
\usepackage{microtype}
\usepackage{xcolor}
\usepackage{listings}
\usepackage{graphicx}
\setcitestyle{numbers,square}

\lstset{
  basicstyle=\ttfamily\small,
  breaklines=true,
  frame=single,
  backgroundcolor=\color{gray!8},
}

\title{Evolving Vision-Language Inference:\\
An AlphaEvolve-Style Code Optimization System\\
for Chart Question Answering}

\author{%
  Qi Huang \\
  The University of Chicago \\
  \texttt{qihuang@uchicago.edu}
}

\begin{document}

\maketitle

\begin{abstract}
We present an AlphaEvolve-style evolutionary code optimization system applied to
Vision-Language Model (VLM) inference for chart understanding on the CharXiv
benchmark. Starting from a naive baseline achieving 38.3\% accuracy, we first
manually optimize the inference pipeline through prompt engineering, structured
output normalization, and image caching, reaching 53.9\% accuracy with a
13$\times$ speed improvement. We then implement an evolutionary system that uses
an LLM (GPT) as a mutation operator over a maintained population of candidate
programs, automatically discovering improvements such as adaptive image
upscaling and unsharp masking, reaching 60.2\% accuracy. Finally, through
systematic per-sample error analysis, we identify and fix five post-processing
bugs (negative sign stripping, text tick mishandling, caret notation breakage,
legend count deduplication, and trend synonym mismatch) that together yield a
further 9.4~pp gain to 69.5\%---a 31.2 percentage point improvement over the
baseline. We also explore test-time augmentation (TTA) and find it net negative
except for a narrow title-verification heuristic. These results show that
careful output analysis and targeted bug fixes can substantially outperform
sophisticated inference techniques when the bottleneck is post-processing rather
than model perception.
\end{abstract}

%%%%%%%%%%%%%%%%%%%%%%%%%%%%%%%%%%%%%%%%%%%%%%%%%%%%%%%%%%%%
\section{Introduction}
%%%%%%%%%%%%%%%%%%%%%%%%%%%%%%%%%%%%%%%%%%%%%%%%%%%%%%%%%%%%

Chart question answering is a challenging multimodal task requiring a model to
read visual artifacts (axes, legends, data lines) and produce exact short-form
answers. The CharXiv benchmark \cite{charxiv} evaluates models on descriptive
and reasoning questions about scientific charts using exact-string comparison,
making answer formatting as important as perceptual accuracy.

AlphaEvolve \cite{alphaevolve} demonstrates that an LLM acting as a mutation
operator inside an evolutionary loop can automatically improve algorithmic code
beyond what human engineering achieves. We apply its core ideas to the VLM
inference pipeline: the ``gene'' being evolved is the Python function
\texttt{vlm\_inference}, and the fitness function is exact-match accuracy on a
128-sample CharXiv development subset.

We fix the answering model to Qwen3-VL-2B \cite{qwen} (Instruct and Thinking
variants) and vary only the surrounding inference code: prompting strategy,
image preprocessing, and post-processing normalization. This clean separation
makes the evolutionary search tractable and interpretable.

%%%%%%%%%%%%%%%%%%%%%%%%%%%%%%%%%%%%%%%%%%%%%%%%%%%%%%%%%%%%
\section{Task Setup}
%%%%%%%%%%%%%%%%%%%%%%%%%%%%%%%%%%%%%%%%%%%%%%%%%%%%%%%%%%%%

\paragraph{Benchmark.}
CharXiv contains 2{,}323 charts from scientific papers with descriptive and
reasoning questions. We use the descriptive validation split and restrict
development/evolution to the first 128 examples as defined in
\texttt{evaluate.py}.

\paragraph{Models.}
We use \texttt{Qwen/Qwen3-VL-2B-Instruct} (direct answer) and
\texttt{Qwen/Qwen3-VL-2B-Thinking} (chain-of-thought reasoning) throughout.
Switching to larger variants is not permitted. Both models are loaded in
\texttt{bfloat16} with \texttt{device\_map="auto"}.

\paragraph{Reproducibility.}
Greedy decoding (\texttt{do\_sample=False}) is enforced in all submitted
scripts. The mutation LLM is OpenAI \texttt{gpt-5.2-codex}; Gemini Flash was
initially attempted but the available project had zero usable
\texttt{generate\_content} quota, so the system was switched to OpenAI.

%%%%%%%%%%%%%%%%%%%%%%%%%%%%%%%%%%%%%%%%%%%%%%%%%%%%%%%%%%%%
\section{Manual Optimization}
%%%%%%%%%%%%%%%%%%%%%%%%%%%%%%%%%%%%%%%%%%%%%%%%%%%%%%%%%%%%

The naive baseline (\texttt{starting\_scripts.py}) loads Qwen3-VL-2B-Instruct,
formats a single-message chat prompt, and generates up to 128 tokens. It
achieves 38.3\% accuracy at 3.48 s/query.

Manual inspection of failure cases showed three dominant error types: (1) the
model produces verbose justifications rather than bare answers; (2) numbers are
formatted with commas or units that the evaluator does not expect; (3) the model
guesses when the chart does not contain evidence, instead of answering ``Not
Applicable''. We address each with targeted changes.

\paragraph{Prompt engineering.}
We replace the default prompt with a structured instruction that explicitly
demands a single-line exact answer, lists the expected format for each question
type (number, Yes/No, \textit{n} by \textit{m}, comma-separated legend labels),
and adds question-conditional clauses (e.g.\ ``if no continuous legend exists,
answer exactly: Not Applicable'').

\paragraph{Output normalization.}
A post-processing layer strips common prefixes (``The answer is'', ``The
x-axis label is''), extracts numbers with a regex, maps layout tokens to the
canonical \textit{n} by \textit{m} form, and collapses ``Not Applicable''
synonyms.

\paragraph{Speed optimizations.}
Image inputs and processed tensors are cached per file path, eliminating
redundant preprocessing for the multiple questions that share the same chart.
\texttt{max\_new\_tokens} is reduced from 128 to 32, and
\texttt{torch.inference\_mode()} replaces \texttt{torch.no\_grad()}.

\paragraph{Thinking model.}
For the Thinking variant, an empty \texttt{<think></think>} block is injected
into the prompt to suppress verbose reasoning chains and force a concise final
answer.

%%%%%%%%%%%%%%%%%%%%%%%%%%%%%%%%%%%%%%%%%%%%%%%%%%%%%%%%%%%%
\section{Evolution System Design}
%%%%%%%%%%%%%%%%%%%%%%%%%%%%%%%%%%%%%%%%%%%%%%%%%%%%%%%%%%%%

Both \texttt{evolve\_instruct.py} and \texttt{evolve\_thinking.py} implement
the core AlphaEvolve components, extended with three mechanisms inspired by the
paper's full architecture.

\paragraph{LLM-based mutation (dual-mode).}
At each generation, the system sends the current parent program along with its
accuracy and speed to \texttt{gpt-5.2-codex}. We implement two mutation modes,
selected with equal probability:
\begin{itemize}
  \item \textbf{Full-block rewrite:} the LLM receives the entire
    \texttt{EVOLVE-BLOCK} and produces a complete replacement, together with a
    randomly chosen mutation strategy from a curated list of eight options.
  \item \textbf{Focused (diff-like) mutation:} we parse the parent block with
    Python's \texttt{ast} module, randomly select one function name, and ask the
    LLM to make a surgical, 1--15 line change to \emph{only} that function while
    keeping everything else identical. This produces smaller, safer mutations
    analogous to AlphaEvolve's diff-based prompting.
\end{itemize}
Both modes include up to two additional population members (including archive
members) as ``inspiration'' in the prompt.

\paragraph{Population management.}
We maintain an island model with 2 islands of 4 programs each (8 total). Every
5 generations the best individual from each island migrates to a random other
island, providing diversity without sacrificing convergence within each island.

\paragraph{Behavioral diversity archive.}
In addition to the island population, we maintain a lightweight MAP-Elites-style
\texttt{FeatureArchive} of up to 20 programs, indexed by the hash of their
per-query correctness vector (a 128-bit behavior signature). Each unique
behavior pattern retains only its highest-scoring representative. This ensures
the system preserves programs that solve \emph{different subsets} of the
benchmark, not just those with the highest aggregate score. Archive members are
sampled as additional ``inspiration'' during mutation, encouraging the LLM to
cross-pollinate solution strategies.

\paragraph{Selection and elitism.}
Tournament selection ($k=3$) within each island chooses the parent for
mutation. A composite fitness score
$f = 0.9 \times \text{accuracy} + 0.1 \times \text{speed\_bonus}$
is used for ranking, keeping accuracy dominant. The globally best program is
checkpointed whenever a new maximum is found.

\paragraph{Cascaded evaluation.}
To avoid wasting compute on broken or obviously weak mutations, evaluation
proceeds in three phases:
\begin{enumerate}
  \item \textbf{Crash check} (1 sample): if the candidate crashes or returns a
    non-string, it is immediately discarded.
  \item \textbf{Pre-screen} (16 samples): the candidate must achieve $\geq 70\%$
    of the parent's accuracy on this subset to proceed. Failed candidates are
    logged and discarded, saving $\sim$87.5\% of evaluation cost.
  \item \textbf{Full evaluation} (all 128 samples): computes final accuracy,
    speed, composite score, and behavior vector.
\end{enumerate}
Phases 1--3 share the same loaded module, so the VLM is only loaded once per
candidate.

\paragraph{Evolution loop.}
Each run performs 40 generations, alternating between islands. Full per-candidate
results (including behavior vectors) are appended under
\texttt{artifacts/evolution\_logs/} for post-hoc analysis.

%%%%%%%%%%%%%%%%%%%%%%%%%%%%%%%%%%%%%%%%%%%%%%%%%%%%%%%%%%%%
\section{Results}
%%%%%%%%%%%%%%%%%%%%%%%%%%%%%%%%%%%%%%%%%%%%%%%%%%%%%%%%%%%%

Table~\ref{tab:results} reports accuracy and average inference time per query
on the 128-sample development subset for all submitted systems.

\begin{table}[h]
\centering
\caption{Accuracy and speed on the 128-sample CharXiv development subset.
Timing is wall-clock per query on the same GPU.}
\label{tab:results}
\begin{tabular}{lrrl}
\toprule
Variant & Accuracy & Time (s/query) & Notes \\
\midrule
\texttt{starting\_scripts} & 0.3828 & 3.484 & naive baseline \\
\texttt{manual\_instruct}  & 0.5391 & 0.269 & +15.6 pp; 13$\times$ faster \\
\texttt{manual\_thinking}  & 0.5078 & 0.275 & thinking model, suppressed reasoning \\
\texttt{evolved\_instruct} & 0.6016 & 0.310 & +6.3 pp over manual (evolution-discovered) \\
\texttt{evolved\_thinking} & 0.5469 & 0.271 & evolved improvement over manual thinking \\
\midrule
\texttt{best\_accuracy}    & \textbf{0.6953} & 0.293 & evolved + post-processing fixes + title verifier \\
\texttt{best\_speed}       & 0.5781 & \textbf{0.265} & no image preprocessing, fast prompt \\
\texttt{best\_overall}     & 0.6797 & 0.310 & evolved + post-processing fixes \\
\bottomrule
\end{tabular}
\end{table}

Table~\ref{tab:progression} shows the incremental accuracy gains that built
the final \texttt{best\_accuracy} result.

\begin{table}[h]
\centering
\caption{Step-by-step accuracy progression for \texttt{best\_accuracy}.}
\label{tab:progression}
\begin{tabular}{lrl}
\toprule
Change & Accuracy & $\Delta$ \\
\midrule
Evolved instruct (baseline for best\_accuracy) & 0.6016 & --- \\
+ Title verifier (3-view TTA, narrow trigger) & 0.6172 & +1.6 pp \\
+ Trend synonym normalization & 0.6250 & +0.8 pp \\
+ Fix negative-sign stripping in \texttt{\_candidates} & 0.6797 & +5.5 pp \\
+ Fix positional-tick text fallback & (combined) & (combined) \\
+ Fix embedded-number regex over-match & (combined) & (combined) \\
+ Preserve caret notation (e.g.\ \texttt{10\^{}-6}) & 0.6953 & +1.6 pp \\
+ Legend count dedup from label lists & (combined) & (combined) \\
\bottomrule
\end{tabular}
\end{table}

\paragraph{Manual vs.\ baseline.}
Manual optimization delivers a 15.6 pp accuracy gain and reduces latency by
roughly 13$\times$, primarily through structured prompting and caching.

\paragraph{Evolution vs.\ manual.}
The evolved instruct variant adds a further 6.3 pp over manual instruct,
reaching 60.2\%. This gain comes from richer image preprocessing (adaptive
upscaling, contrast enhancement, unsharp masking) and tighter prompt
constraints discovered by the mutation LLM. Notably, these image-level changes
were not part of the manual baseline and would have been non-obvious to add by
hand.

\paragraph{Post-processing bug fixes.}
The largest accuracy leap (+9.4~pp, from 0.6016 to 0.6953) came not from
model-level changes but from systematic per-sample error analysis of raw model
outputs vs.\ post-processed answers. We identified five bugs:
\begin{enumerate}
  \item \textbf{Negative sign stripping:} \texttt{line.strip(" -:")} in
    \texttt{\_candidates()} removed the leading minus from values like
    \texttt{-1.00}, causing mismatches on tick-value questions.
  \item \textbf{Positional tick text fallback:} tick labels such as ``SSM'' or
    ``Heter-unaware'' were always routed to the numeric extraction branch,
    which returned ``Not Applicable'' instead of the text answer.
  \item \textbf{Embedded-number regex over-match:} a label like
    ``ResNet-1'' caused the number regex to extract ``1'' instead of returning
    the full text string.
  \item \textbf{Caret scientific notation:} \texttt{10\^{}-6} was split by
    the number regex into separate tokens; we now detect the caret pattern
    and preserve it as-is.
  \item \textbf{Legend count dedup:} when asked ``how many discrete labels,''
    the model sometimes outputs a comma-separated label list instead of a
    count; we now split, deduplicate, and count.
\end{enumerate}

\paragraph{Instruct vs.\ Thinking.}
The Instruct model consistently outperforms the Thinking model across both
manual and evolved variants. The Thinking model's chain-of-thought mechanism is
a liability for exact-match short-answer tasks: even with reasoning suppressed,
it is more sensitive to prompt phrasing and occasionally produces longer
responses that fail normalization.

\paragraph{Test-time augmentation (TTA).}
We explored multi-view TTA---generating answers from multiple image
enhancements (sharpened, grayscale, high-contrast) and aggregating via majority
vote. \textbf{Global TTA was net negative} (0.6016$\to$0.5625): the additional
views introduced more noise than signal across most question types. We then
narrowed the scope to a \textbf{title verification heuristic}: if the base
answer to a ``what is its title'' question matches a panel-marker pattern
(e.g.\ ``(b)'', ``(d)''), we run 3 views with a strict probe prompt and
require unanimous agreement to override. This rescued 2 samples with zero
injuries (+1.6~pp).

We also tested consistency-check mediators for legend count, legend label, and
colorbar questions. All were unsafe: the model's meta-judgments (``does this
chart have a legend?'') were less reliable than its direct answers.

\paragraph{Best deliverables.}
Based on the above:
\begin{itemize}
  \item \textbf{Best accuracy:} \texttt{best\_accuracy} (0.6953) --- evolved
        base + all post-processing fixes + title verifier.
  \item \textbf{Best speed:} \texttt{best\_speed} (0.269 s/query) --- no image
        preprocessing, lightweight prompt, same post-processing fixes.
  \item \textbf{Best overall:} \texttt{best\_overall} --- evolved image
        preprocessing + all post-processing fixes, no title verifier (to avoid
        the 3$\times$ latency penalty on title questions).
\end{itemize}

%%%%%%%%%%%%%%%%%%%%%%%%%%%%%%%%%%%%%%%%%%%%%%%%%%%%%%%%%%%%
\section{Analysis}
%%%%%%%%%%%%%%%%%%%%%%%%%%%%%%%%%%%%%%%%%%%%%%%%%%%%%%%%%%%%

\paragraph{What worked.}
Exact-match-oriented prompting had the highest single-change impact in the
manual phase. Many baseline failures were pure formatting failures: the model
knew the correct value but wrapped it in a sentence. Structured output
normalization provided additional gains at essentially zero computational cost.
Seeding evolution from an already-optimized manual baseline (rather than the
naive starter) was critical: early generations can refine rather than recover
basic competence.

The most impactful technique overall was \textbf{per-sample error analysis},
comparing raw model output against post-processed answers. This revealed that
the model was often correct but the normalization code was discarding or
corrupting the answer. Five targeted bug fixes yielded +9.4~pp at zero
computational cost.

\paragraph{What did not work.}
Global test-time augmentation (multi-view voting) was net negative ($-$3.9~pp).
The 2B model's outputs were highly consistent across image variants, so
majority voting rarely rescued errors but did introduce new ones when a minority
view happened to be wrong.

Consistency-check mediators (asking the model meta-questions like ``does this
chart have a legend?'' to gate answers) were also unsafe. The model's
meta-judgments proved less reliable than its direct task answers.

A non-trivial fraction of GPT mutations (8/37 for instruct, 1/32 for thinking)
produced syntactically valid but semantically broken code (zero accuracy at
runtime). More robust sanity-checking (e.g.\ a one-sample pre-screen) could
save significant evaluation time.

\paragraph{Remaining errors.}
Of the 39 remaining errors (out of 128), the dominant categories are:
\begin{itemize}
  \item \textbf{Colorbar existence confusion} (13 errors): the model
    confidently reads tick values from charts that have no colorbar, or
    reads values from the wrong chart region.
  \item \textbf{Numeric tick counting bias} (7 errors): systematic
    over-counting when asked ``how many ticks across all axes.''
  \item \textbf{OCR precision} (remainder): the model reads
    visually similar numbers incorrectly (e.g.\ ``0.7'' vs ``0.8'').
\end{itemize}
These are model perception limits, not fixable through post-processing.

\paragraph{Future directions.}
Increasing the number of generations, maintaining a larger population, or using
a more powerful mutation model could push accuracy further. Adding cascaded
evaluation (fast pre-screen on a small sample subset) and behavioral diversity
maintenance (MAP-Elites-style feature archives) to the evolution system would
improve both efficiency and exploration.

%%%%%%%%%%%%%%%%%%%%%%%%%%%%%%%%%%%%%%%%%%%%%%%%%%%%%%%%%%%%
\section{Conclusion}
%%%%%%%%%%%%%%%%%%%%%%%%%%%%%%%%%%%%%%%%%%%%%%%%%%%%%%%%%%%%

We implemented an AlphaEvolve-style evolutionary code optimization system for
VLM chart inference and showed that it improves accuracy by 21.9 pp over the
naive baseline and 6.3 pp over careful manual optimization. Subsequent
per-sample error analysis and targeted post-processing bug fixes pushed accuracy
to 69.5\%---a 31.2 pp total improvement. The results confirm that automated
code evolution, seeded with human insight, is a practically useful starting
point, but that systematic output analysis and targeted engineering remain
essential for maximizing exact-match accuracy when the model weights are fixed.

%%%%%%%%%%%%%%%%%%%%%%%%%%%%%%%%%%%%%%%%%%%%%%%%%%%%%%%%%%%%
\bibliographystyle{plain}
\begin{thebibliography}{9}

\bibitem{alphaevolve}
Google DeepMind.
\newblock AlphaEvolve: A Gemini-powered coding agent for designing advanced
algorithms, 2025.
\newblock \url{https://deepmind.google/discover/blog/alphaevolve-a-gemini-powered-coding-agent-for-designing-advanced-algorithms/}

\bibitem{charxiv}
Z.~Wang et al.
\newblock CharXiv: Charting Gaps in Realistic Chart Understanding in Multimodal
LLMs.
\newblock \textit{NeurIPS}, 2024.

\bibitem{qwen}
Qwen Team.
\newblock Qwen3 Technical Report, 2025.
\newblock \url{https://huggingface.co/Qwen/Qwen3-VL-2B-Instruct}

\end{thebibliography}

\end{document}
